\documentclass[12pt]{article}
\usepackage{psfrag}
\usepackage{epsf,epsfig,amsfonts,amssymb}
\setlength{\textwidth}{165 mm}
\setlength{\oddsidemargin}{-0.05in}
\setlength{\textheight}{9.5in}
\setlength{\topmargin}{-0.7in}
\pagestyle{plain}
\hoffset=5mm
\usepackage[centertags]{amsmath}    
\usepackage{graphicx}
\usepackage{amsfonts}
\usepackage{amssymb}
\usepackage{amsmath}
\usepackage{subfigure} 
\usepackage{array}
\usepackage{float}



\begin{document}

\noindent
{\it Proceedings of the First UK Conference on Inverse Problems}\\
(eds. D. Lesnic and N. Kinash), Leeds University Press, pp. ?-?, 2027\\

\noindent
\\

\begin{center}
\Large{\bf Inverse interface determination using a level-set approach}
\end{center}

\begin{center}
Initials. Surname$^{1}$ and Initials. Surname$^{2}$\\
$^{1}${\it Department ..., University..., Country. E-mail: }\\
$^{2}${\it Department ..., University..., Country. E-mail: }\\
\end{center}

\noindent
\\

\noindent
{\bf Abstract.} The paper should have maximum 10 pages.

\section{Introduction} \label{Section1}
Permeability spans several order of magnitudes over different length scales and 
dynamic data from wells provide important information for permeability 
estimation.\\
\indent
This classical but fundamental inverse problem of groundwater flows, as named in the review paper by \cite{Yeh}, with obvious importance to the oil industry, is naturally suited 
to a zonation approach in which the permeability of each zone of the inhomogeneous aquifer is unknwon 
and has to be inferred from piezometric head 
measurements. 

\section{Mathematical formulation} \label{Section2}
Consider the steady-state reaction-diffusion equation
\begin{eqnarray}
- \nabla \cdot (a({\bf x}) \nabla u) + C({\bf x}) u=Q({\bf x}), \quad {\bf x} \in \Omega:=(-1,1)^{2}. 
\label{eq:1}
\end{eqnarray}

\section{Conclusions} \label{Section5}
A simple, easy-to-use and fast regularised level-set approach has been devised to 
find the interfaces between materials.  \\

\noindent
{\bf Acknowledgements}\\
Some financial support received from the University of Leeds is acknowledged.\\

\noindent
Figures should be black and white, i.e. no colour. Use different line styles or markers. The numbering and lettering on the axes should be big.

\begin{figure}[H]
\begin{center}
\includegraphics[scale=0.5]{Objective_func_case1_r1.eps}
\end{center}
\caption{Logarithm of the objective function.}
\label{Fig1}
\end{figure}

Please use consistently the reference style below.


\begin{thebibliography}{10}
\bibitem{Ang}{W.T. Ang and H. Fan} (2003) A hypersingular boundary integral formulation for heat conduction across an imperfect interface, In: {\it Advances in Boundary Element Techniques IV}, (eds. R. Gallego and M.H. Aliabadi), Queen Mary, University of London, UK, pp.267-274.

\bibitem{Berre1}{I. Berre, M. Lien and T. Mannseth} (2008) A multi-level strategy for regularized parameter identification 
using nonlinear reparametrisation 
with sample application for reservoir characterization, {\it Journal of Physics: Conference Series}, {\bf 135}, 
Article 012018, (8 pages).

\bibitem{Suetin}{P.K. Suetin} (1999) {\it Orthogonal Polynomials in Two Variables}, Gordon and Breach, London.

\bibitem{Yeh}{W.W.-G. Yeh} (1986) Review of parameter identification procedures in groundwater hydrology: The inverse problem, {\it Water Resources Research}, {\bf 22}, 95-108.
\end{thebibliography}

\end{document}

